\newif\ifblind\blindfalse
\newif\iflipics\lipicsfalse
\usepackage{amsmath,amsthm}
\usepackage{fullpage}
\usepackage{thm-restate}

\usepackage[pagebackref]{hyperref}
\usepackage[capitalize,nameinlink]{cleveref}
\usepackage[T1]{fontenc}
\usepackage[utf8]{inputenc}
\usepackage{amsfonts,amssymb,bm,bbm}
\usepackage{mathtools}
\usepackage{braket}
\usepackage{physics}
\usepackage{thmtools}
\usepackage{graphicx}
\usepackage{enumitem}
\usepackage{xcolor}
\usepackage{tikz}
\usetikzlibrary{matrix}
\usepackage{xspace}
\usepackage{graphicx}
\usepackage{subcaption}
\usepackage{physics}
\usepackage[breakable]{tcolorbox}
\usepackage{crossreftools}

\makeatletter
\newcommand{\tcbline}{%
  \begin{tikzpicture}
  \path[use as bounding box] (0,0) -- (\linewidth,0);
  \draw[color=black!20!white,dashed,dash phase=2pt]
        (0-\kvtcb@leftlower-\kvtcb@boxsep,0)--
        (\linewidth+\kvtcb@rightlower+\kvtcb@boxsep,0);
  \end{tikzpicture}%
  }
\makeatother

\makeatletter
\def\@footnotecolor{red}
\define@key{Hyp}{footnotecolor}{%
 \HyColor@HyperrefColor{#1}\@footnotecolor%
}
\def\@footnotemark{%
    \leavevmode
    \ifhmode\edef\@x@sf{\the\spacefactor}\nobreak\fi
    \stepcounter{Hfootnote}%
    \global\let\Hy@saved@currentHref\@currentHref
    \hyper@makecurrent{Hfootnote}%
    \global\let\Hy@footnote@currentHref\@currentHref
    \global\let\@currentHref\Hy@saved@currentHref
    \hyper@linkstart{footnote}{\Hy@footnote@currentHref}%
    \@makefnmark
    \hyper@linkend
    \ifhmode\spacefactor\@x@sf\fi
    \relax
  }%
\makeatother

\hypersetup{
  linkcolor  = blue!80!black,
  citecolor  = green!50!black,
  footnotecolor = red!90!black,
  colorlinks = true,
}

\makeatletter
\newcommand{\subalign}[1]{%
  \vcenter{%
    \Let@ \restore@math@cr \default@tag
    \baselineskip\fontdimen10 \scriptfont\tw@
    \advance\baselineskip\fontdimen12 \scriptfont\tw@
    \lineskip\thr@@\fontdimen8 \scriptfont\thr@@
    \lineskiplimit\lineskip
    \ialign{\hfil$\m@th\scriptstyle##$&$\m@th\scriptstyle{}##$\hfil\crcr
      #1\crcr
    }%
  }%
}
\makeatother

\newlist{steps}{enumerate}{10}
\setlist[steps]{label=Step \arabic*:, ref=\arabic*}
\crefname{stepsi}{Step}{Steps}

\makeatletter
\newcommand{\envfootnote}{%
\addtocounter{footnote}{+1}%
\addtocounter{Hfootnote}{+1}%
\global\let\Hy@saved@currentHref\@currentHref%
\hyper@makecurrent{Hfootnote}%
\global\let\Hy@footnote@currentHref\@currentHref%
\global\let\@currentHref\Hy@saved@currentHref%
}
\makeatother

\iflipics
\else
\newtcolorbox[
  auto counter,
  number within=section,
  number freestyle={\noexpand\thesection.\noexpand\arabic{\tcbcounter}},
  crefname={Protocol}{Protocols}]%
{protocol}[2][]{
  breakable,
  colback=black!3!white,
  colframe=black!20!white,
  coltitle=black,
  code={\def\mytitle{\normalfont #2}},
  fonttitle=\bfseries,
  title=Protocol \thetcbcounter: \mytitle, #1,
  beforeafter skip=\baselineskip}

\newtcolorbox[
  auto counter,
  number within=section,
  number freestyle={\noexpand\thesection.\noexpand\arabic{\tcbcounter}},
  crefname={Algorithm}{Algorithms}]%
{algorithm}[2][]{
  breakable,
  colback=black!3!white,
  colframe=black!20!white,
  coltitle=black,
  code={\def\mytitle{\normalfont #2}},
  fonttitle=\bfseries,
  title=Algorithm \thetcbcounter: \mytitle, #1,
  beforeafter skip=\baselineskip}

\creflabelformat{equation}{#2#1#3}
\newtheorem{theorem}{Theorem}[section]
\newtheorem{lemma}[theorem]{Lemma}
\newtheorem{corollary}[theorem]{Corollary}
\counterwithin{corollary}{section}
\newtheorem{claim}[theorem]{Claim}
\newtheorem{proposition}[theorem]{Proposition}

\theoremstyle{definition}
\newtheorem{remark}[theorem]{Remark}
\newtheorem{definition}[theorem]{Definition}
\fi
\newtheorem{openproblem}{Open problem}
\newcommand{\zksip}{\textrm{zkSIP}}
\newcommand{\eat}[1]{}
\DeclareMathOperator{\view}{View}
\newcommand{\E}{\mathbb E}
\newcommand{\N}{\mathbb N}
\newcommand{\Z}{\mathbb Z}
\newcommand{\field}{\mathbb{F}}
\newcommand{\Field}{\mathbb{H}}
\newcommand{\R}{\mathbb R}
\renewcommand{\P}{\mathbb P}
\renewcommand{\set}[1]{\left\{#1\right\}}
\newcommand{\poly}{\operatorname{poly}}
\newcommand{\polylog}{\operatorname{polylog}}
\newcommand{\eps}{\varepsilon}
\newcommand{\bitset}{\{0,1\}}
\newcommand{\indexproblem}{\normalfont\textsc{index}\xspace}
\newcommand{\searchindex}{\normalfont\textsc{search-index}\xspace}
\newcommand{\decisionindex}{\normalfont\textsc{decision-index}\xspace}
\newcommand{\pair}{\normalfont\textsc{pair}\xspace}
\newcommand{\reconstruct}{\normalfont\textsc{reconstruct}\xspace}
\newcommand{\pep}{\normalfont\textsc{pep}\xspace}
\newcommand{\pepprotocol}{\textsf{pep}\xspace}
\newcommand{\sumcheck}{\normalfont\textsc{sumcheck}\xspace}
\newcommand{\frequencymoment}{\normalfont\textsc{frequency-moment}\xspace}
\newcommand{\pointquery}{\normalfont\textsc{point-query}\xspace}
\newcommand{\rangecount}{\normalfont\textsc{range-count}\xspace}
\newcommand{\selection}{\normalfont\textsc{selection}\xspace}
\newcommand{\innerproductproblem}{\normalfont\textsc{inner-product}\xspace}
\newcommand{\median}{\normalfont\textsc{median}\xspace}
\newcommand{\vcommitlength}{v}
\newcommand{\vcommitlengthvalue}{n^{2\log\log n}\log n}
\newcommand{\vspacevalue}{O(\log^2 n)}
\newcommand{\vcommitstring}{z}
\newcommand{\pcommitlength}{p}
\newcommand{\pcommitlengthvalue}{n^{2.01 \log n}}
\newcommand{\pcommitstring}{y}
\newcommand{\dimension}{m}
\newcommand{\dimensionvalue}{\frac{\log n}{\log\log n}}
\newcommand{\degree}{d}
\newcommand{\degreevalue}{\log n}
\newcommand{\Dimension}{{m'}}
\newcommand{\Dimensionvalue}{O(\log^2 n/\log\log n)}
\newcommand{\Degree}{{d'}}
\newcommand{\Degreevalue}{O(\log^2 n)}
\newcommand{\fieldsize}{q}
\newcommand{\fieldsizevalue}{\log^2 n}
\newcommand{\Fieldsize}{{q'}}
\newcommand{\Fieldsizevalue}{\log^4 n}
\newcommand{\evaluationpoint}{\bm{\fe2}}
\newcommand{\fingerprinttuple}{\bm{\rho}}
\newcommand{\Fingerprinttuple}{\bm{\sigma}}
\renewcommand{\line}{L}
\newcommand{\snapshot}{b}
\newcommand{\whiteboxoracle}{\mathcal{W}}
\newcommand{\cubeside}{H}
\newcommand{\cubesidelength}{h}
\newcommand{\ball}{\mathcal B}
\newcommand{\alphabet}{\Gamma}
\newcommand{\alphabetsize}{\gamma}
\newcommand{\chf}{\mathbbm 1}

\def\symbol#1{
    \ifnum#1=1
        \alpha
    \else
    \ifnum#1=2
        \beta
    \else
    \ifnum#1=3
        \gamma
    \fi\fi\fi
}

\def\fe#1{
    \ifnum#1=1
        \alpha
    \else
    \ifnum#1=2
        \beta
    \else
    \ifnum#1=3
        \gamma
    \fi\fi\fi
}

\def\rfe#1{
    \ifnum#1=1
        \rho
    \else
    \ifnum#1=2
        \sigma
    \else
    \fi\fi
}
\graphicspath{{items/}}
